\section{A direct majorant of the complete Pell radical}
\label{sec:directmajorant}

Another new construction improves a different measure of simplicity:
the total degree with nine positive unknowns.  It retains package~71's
source encoding and replaces its final square-root bound.  The resulting
polynomial has degree $1\,158\,228$, compared with $1\,379\,730$ for that
package, and its explicit circuit has $1\,753$ arithmetic operations,
compared with $478$.  Thus this is a degree--circuit-size tradeoff.
The improvement is approximately $16.05\%$ in degree.  It neither lowers
the number of unknowns nor asserts a global minimum.

\subsection{Why approximate the complete radical?}

The previous finalizer constructs a polynomial upper bound for
$\sqrt{DF}\,G$ and then multiplies it by $H$.  Errors in that bound are
therefore multiplied by a high-degree polynomial.  We instead construct
a polynomial upper bound for the complete radical
\[
 \sqrt{DF\bigl((G^2-1)H^2+1\bigr)}.
\]
The discovery calculation uses a formal expansion at $A=\infty$.
The proof of the resulting bound consists entirely of finite integer
polynomial identities and positive coefficients.

In this subsection $A,C,i,Z$ are independent variables.  Define
\begin{align*}
 J&=4A-5,& D&=(A^2-1)C^2+1,\\
 e&=iJC,& E&=eDC,\\
 F&=(A^2-1)E^2+1,& L_0&=CD+1-A,\\
 G&=A+FL_0,& H&=C+FZ,\\
 I&=(G^2-1)H^2+1,& Q_*&=DFI.
\end{align*}
As a polynomial in $A$, $Q_*$ has degree $46$ and leading coefficient
\[
 (1024C^{24}i^5Z)^2.
\]
Write $Q_*=\sum_{k=0}^{46}q_k A^k$.  Define $t_{23},\ldots,t_0$ in
$\mathbb Q(C,i,Z)$ by the finite recurrence
\begin{align}
 t_{23}&=1024C^{24}i^5Z,\nonumber\\
 t_{23-k}
 &=\frac{q_{46-k}-\sum_{\ell=1}^{k-1}
                 t_{23-\ell}t_{23-k+\ell}}{2t_{23}}
       \qquad(1\leq k\leq23).                 \label{eq:majorantrecurrence}
\end{align}
This is the polynomial part of the formal square root with positive
leading term.  The recurrence is a finite specification of coefficients;
no division by an unknown remains in the final equation.
Set
\begin{equation}
 \mathcal P(A,C,i,Z)=65536C^7\sum_{r=0}^{23}t_r A^r,
 \qquad \sigma=65536C^7.                       \label{eq:majorantP}
\end{equation}
Exact coefficient arithmetic proves
$\mathcal P\in\mathbb Z[A,C,i,Z]$.  It has $654$ nonzero terms, all of
which are supplied in \texttt{nine\_positive/data/radical.json}.
The same file gives every coefficient of
\begin{equation}
 \mathcal D=\mathcal P^2-\sigma^2Q_*.
                                                    \label{eq:localdefect}
\end{equation}
The coordinate degrees of these two polynomials are
\[
 \deg_{(A,C,i,Z)}\mathcal P=(23,31,5,1),\qquad
 \deg_{(A,C,i,Z)}\mathcal D=(22,42,10,2).
\]
There are $3781$ nonzero defect terms.  Denominators in the formal
coefficient recurrence cancel after the displayed factor $65536C^7$;
there are no negative exponents in either delivered table.

\subsection{A uniform inequality certified by positive coefficients}

The decisive certificates are the following exact identities, represented
by their complete coefficient lists:
\begin{enumerate}
 \item $\mathcal P(X+3,Y+1,i,Z)$ has $1816$ nonzero coefficients,
       all strictly positive.
 \item $\mathcal D(X+3,Y+1,U+1,V+1)$ has $32\,637$ nonzero
       coefficients, all strictly positive.  Its constant coefficient is
       \[
          295870595720775501924817.
       \]
\end{enumerate}
The lists are in \texttt{certificates/P\_shifted.json} and
\texttt{certificates/defect\_shifted.json} within the nine-positive
component.  Regeneration by rational coefficient arithmetic and an
independent implementation using only integer sparse polynomials agree.
These are finite algebraic certificates for statements over the whole
domain, rather than sampled inequalities.

\begin{lemma}[Complete-radical majorant]
For all integers $A\geq3$ and $C,i,Z\geq1$,
\begin{equation}
 \mathcal P(A,C,i,Z)>65536C^7\sqrt{DFI}>0.
                                                    \label{eq:majorantineq}
\end{equation}
\end{lemma}
\begin{proof}
Here $J,D,E,F$ are positive.  Since
$CD+1-A\geq A^2-A+1>0$, we have $G>A\geq3$; also $H,I>0$.
The first certificate evaluated at $X=A-3$, $Y=C-1$ gives
$\mathcal P>0$.  All its monomials are nonnegative, and its terms of
$X$- and $Y$-degree zero are positive when $i,Z\geq1$.
The second certificate evaluated at
$(X,Y,U,V)=(A-3,C-1,i-1,Z-1)$ gives $\mathcal D>0$.
Strictness holds even on every boundary because its constant is positive.
Equation~\eqref{eq:localdefect}, together with $\mathcal P>0$, now gives
\eqref{eq:majorantineq}.
\end{proof}

The use of a formal expansion in the discovery process therefore does
not introduce an asymptotic qualification into the theorem.  The bound
holds before any solution of the final equation has been assumed.

\subsection{The fixed equation and its nine unknowns}

All nine unknowns
\[
 b,g,h,s,w,C,i,j,n\in\mathbb Z_{>0}
\]
retain the domain and meaning in the copied source-71 prefix.  The source
polynomials $R,N,t$ are the fixed source-65 encoding polynomials; their
complete weights, target positions and coefficient kernels are included
in the component's data directory.  The whole original-31 native
coefficient tuple is supplied by literal payloads and a finite generator.
No coefficient is left as a free program parameter.

For clarity, the ratio registers needed below are
\begin{align*}
 U&=N^2w,&Y&=N^2s,& M&=RY,\\
 P_0&=2M^2U,&K&=R+1+h(P_0-1),&A&=M(U+1),\\
 B_1&=2R+1,&W_1&=bw,&J&=4A-5,\\
 Q_0&=3W_1C-2(W_1^2-1),&&&&\\[-6pt]
 m&=K^2(4t-g)-16t(C-KY)^2,&&&&\\
 \Omega&=Jm,&E_0&=J(n+W_1^2)+Q_0,&V_0&=P_0K,\\
 a_0&=(P_0^2-1)K^2+1,& Z&=B_1+(j-1)C.&&
\end{align*}
In the local polynomials above, substitute this $A$, this $C$, the free
$i$, and this $Z$.  In particular $H=C+FZ$ is exactly the inherited
Lucas-certificate register.  Define
\begin{equation}
 V_1=\mathcal P(A,C,i,Z),\qquad
 a_1=(65536C^7)^2DFI,\qquad
 \Theta=E_0-\Omega(V_0+V_1).                 \label{eq:newnineinputs}
\end{equation}
The new equation is
\begin{equation}
 \boxed{\quad
 F_9=(\Theta^2-\Omega^2(a_0+a_1))^2
                       -4\Omega^4a_0a_1=0.\quad}   \label{eq:newnine}
\end{equation}
Every named register is a polynomial abbreviation, not an additional
unknown.  The division-free DAG \texttt{equation9.json} is the complete
literal instance.  The source bounds $R,N\geq8$ imply $A\geq3$ on every
ambient positive input, and $C,i,Z\geq1$ is immediate.  Consequently the
majorant lemma applies without a circular appeal to the final equation.

\subsection{Why the replacement preserves the ratio theorem}

The original first majorant satisfies $V_0>\sqrt{a_0}>0$, since
$V_0^2-a_0=K^2-1>0$.  The lemma gives
$V_1>\sqrt{a_1}>0$.  Moreover
\[
 E_0=Jn+(J-2)W_1^2+3W_1C+2>0.
\]
Thus a zero of~\eqref{eq:newnine} cannot have $m=0$: that would give
$F_9=E_0^4>0$.  Factoring the two-root norm over the reals gives signed
roots $r_0^2=a_0$, $r_1^2=a_1$ with
\begin{equation}
 \frac{E_0}{J}=m(V_0+V_1+r_0+r_1).             \label{eq:marginorientation}
\end{equation}
The strict majorants make the parenthesis positive for every sign choice,
so $m>0$ follows before any root-integrality argument.  The unchanged
margin argument now recovers the required free-carrier bounds and makes
$DFI>a_0$, hence $a_1>a_0$.

The sum $r_0+r_1=\Theta/\Omega$ is rational and nonzero.  The identity
\[
 2(r_0+r_1)r_0=(r_0+r_1)^2+a_0-a_1
\]
therefore makes both roots rational, and an integer with a rational square
root has an integral square root.  Although the scale
$65536C^7$ depends on the current positive input, it is a nonzero integer.
Consequently $r_1/(65536C^7)$ is a rational square root of $DFI$ and is
also integral.  No fixed-scale assumption is needed.  Equation
\eqref{eq:marginorientation} then recovers $J\mid Q_0$ exactly as before.
All the inherited ratio conditions, and hence the source zero, follow.

Conversely, a source zero supplies the original ratio witnesses and
integral roots.  Choose the positive roots
$r_0=\sqrt{a_0}$ and $r_1=65536C^7\sqrt{DFI}$, and take
\[
 n=m(V_0+V_1+r_0+r_1)-W_1^2-Q_0/J.
\]
This is an integer.  The unchanged estimate $W_1^2+Q_0/J<2D$, together
with $F>D$, $I\geq1$ and $V_1>65536\sqrt{DFI}>2D$, shows $n>0$.
One norm factor vanishes.  This proves the same source-zero equivalence
as source~71.  Composing with the retained native-source and syntactic
refutation calibration, and negating, gives the ordinary theorem
formalizable in ZFC:
\[
 \mathsf{ZFC}\vdash
 \bigl(\text{$F_9$ has no positive integer zero}\bigr)
       \ \longleftrightarrow\ \operatorname{Con}(\mathsf{ZFC}).
\]
The new component supplies no additional end-to-end Lean or
object-language ZFC certificate; the inherited calibration boundary is
stated explicitly in its proof and provenance files.

\subsection{Exact degree and the change of dominant term}

Let $r=7963$, $\nu=2115$, $W=116$ be the degrees of $R$, $N$ and the
maximum source weight, and put $a=r+4\nu+2=16425$.  Then
\begin{align*}
 \deg\Omega&=5a+2W+3=82360,&\deg V_0&=3a+r+1=57239,\\
 \deg\delta_0&=3a+r+2=57240,&
 \delta_0&=V_0^2-a_0=K^2-1.
\end{align*}
The unique highest weighted-degree term of $\mathcal P$ is
$2^{26}A^{23}C^{31}i^5Z$, giving
\[
 \deg V_1=23a+r+36=385774.
\]
The corresponding unique term of its defect is
\[
 141664530202624\,A^{22}C^{42}i^{10}Z^2,
\]
so $\delta_1=V_1^2-a_1$ has degree $22a+2r+52=377328$.
To resolve the cancellations in the outer norm, set
\[
 B_*=E_0^2-2E_0\Omega(V_0+V_1)+\Omega^2(\delta_0+\delta_1).
\]
It has degree $542048$, and the exact identity
\begin{equation}
 F_9=B_*^2+4\Omega^2V_0V_1B_*
                 +4\Omega^4(\delta_0V_1^2+a_0\delta_1)       \label{eq:resolvednorm}
\end{equation}
shows the following four degrees:
\begin{center}
\begin{tabular}{lr}
$B_*^2$ & $1\,084\,096$\\
$4\Omega^2V_0V_1B_*$ & $1\,149\,781$\\
$4\Omega^4\delta_0V_1^2$ & $1\,158\,228$\\
$4\Omega^4a_0\delta_1$ & $821\,246$
\end{tabular}
\end{center}
The third term is uniquely largest.  Its leading coefficient is positive,
so the exact degree is
\[
 \boxed{\deg F_9=72r+276\nu+8W+224=1\,158\,228.}
\]
The dominant term is different from that of the previous finalizer;
merely substituting the new defect degree into its old formula would be
incorrect.  If $\rho,\sigma_N,\tau$ are the positive diagonal leading
coefficients of $R,N,t$, respectively, the new leading coefficient is
$2^{90}\tau^4\rho^{72}\sigma_N^{276}>0$.  For the actual source constants
$\alpha=2^{2134939}$ and $c=2^{17079501}$ this becomes
\[
 2^{162}\alpha^{578828}c^{276}>0.
\]
Thus specialization to the actual native coefficients cannot remove the
leading term.

The final circuit contains $923$ multiplications, $810$ additions and
$20$ subtractions.  This count includes every multiplication by a fixed
integer and counts reuse once.  The independent verifier reconstructs
all $654$ majorant coefficients directly from the final DAG, checks the
whole finalizer as a polynomial identity, and verifies $45$ retained
named source expressions against source~71.  Both independent checkers
use only the Python standard library and pass with and without Python's
optimization flag.
